%% =====================================================================
%%  B3-T-12-03 -- what B3 contributes to "Friends and Enemies"
%%  -------------------------------------------------------------------
%%  LAYER A: APPEND-ONLY. Written once at the entry's write, then left
%%  alone. If the source has more to say later, add a bullet -- do not
%%  rewrite. Material found ONLY here goes in the `unique` slot with
%%  @unique set to yes, which makes it undeletable (rule R10).
%% =====================================================================
%% @kind: source
%% @id: B3-T-12-03
%% @book: B3
%% @topic: T-12-03
%% @locator: Law 2, pp. 13--18
%% @status: distilled
%% @unique: no
%% @integrated: yes
%% @updated: 2026-09-19

\dossier{B3}{T-12-03}{\bookshort{B3}}{Law 2, pp. 13--18}

\begin{distilled}
  \item Be wary of friends: they betray more quickly, being easily aroused to envy; if you have no enemies, find a way to make them.
  \item Never expect gratitude from a friend and you will be pleasantly surprised; the problem with hiring friends is that it limits power --- the friend is rarely the most able.
  \item The jaws of ingratitude image: you would not put a finger in a lion's mouth, but with friends you drop all caution, and hired friends eat you alive with ingratitude.
  \item Keep enemies as enemies rather than transforming them: the reconciled enemy has everything to prove and no illusions to outgrow (Mao's use of the Nationalist--Japanese conflict).
\end{distilled}
